\chapter{AI Model Evidence}
This appendix records the AI prediction evidence computed from the included dataset and saved artifacts. The values describe the repository's synthetic dataset and should not be interpreted as validation on real HR data.

\section{Dataset Fields}
\begin{longtable}{p{0.28\textwidth}p{0.18\textwidth}p{0.44\textwidth}}
\caption{AI dataset and prediction fields.}\label{tab:ai-fields}\\
\toprule
Field & Used by model & Meaning\\
\midrule
\endfirsthead
\toprule
Field & Used by model & Meaning\\
\midrule
\endhead
\texttt{kpi\_score} & Yes & Numeric KPI achievement score between 0 and 100.\\
\texttt{goal\_completion\_percent} & Yes & Percentage of goals completed.\\
\texttt{checkin\_count} & Yes & Number of submitted check-ins, bounded by the API between 0 and 20.\\
\texttt{avg\_checkin\_progress} & Yes & Average progress reported in check-ins.\\
\texttt{feedback\_count} & Yes & Number of feedback items.\\
\texttt{positive\_feedback\_ratio} & Yes & Ratio between 0 and 1.\\
\texttt{task\_completion\_percent} & Yes & Percentage of assigned tasks completed.\\
\texttt{tasks\_on\_time\_percent} & Yes & Percentage of tasks completed on time.\\
\texttt{overall\_score} & No & Deterministic score excluded from training to avoid leakage.\\
\texttt{rating} & Target & Four-class performance rating target.\\
\texttt{promotion\_ready} & Target & Binary promotion-readiness target.\\
\texttt{strengths}, \texttt{weaknesses}, \texttt{review\_summary}, \texttt{suggestions} & No & Generated text fields used for realistic API outputs.\\
\bottomrule
\end{longtable}

\section{Dataset Quality Checks}
\begin{table}[H]\centering\small
\caption{Dataset checks computed locally.}
\label{tab:ai-quality-checks}
\begin{tabular}{lr}
\toprule
Check & Value\\
\midrule
Rows & 10,000\\
Missing values & 0\\
Duplicate rows & 0\\
Rating classes & 4\\
Prediction input features & 8\\
\bottomrule
\end{tabular}
\end{table}

\section{Class Distribution}
\begin{table}[H]\centering\small
\caption{Rating class distribution in the synthetic dataset.}
\label{tab:ai-class-distribution}
\begin{tabular}{lr}
\toprule
Rating & Count\\
\midrule
\texttt{exceptional} & 2,238\\
\texttt{exceeds\_expectations} & 2,722\\
\texttt{meets\_expectations} & 2,440\\
\texttt{needs\_improvement} & 2,600\\
\bottomrule
\end{tabular}
\end{table}

\section{Promotion Readiness by Rating}
\begin{table}[H]\centering\small
\caption{Promotion-readiness rate by rating in the synthetic data.}
\label{tab:promotion-rate}
\begin{tabular}{lr}
\toprule
Rating & Promotion-ready rate\\
\midrule
\texttt{exceptional} & 0.7962\\
\texttt{exceeds\_expectations} & 0.5103\\
\texttt{meets\_expectations} & 0.2102\\
\texttt{needs\_improvement} & 0.0246\\
\bottomrule
\end{tabular}
\end{table}

\section{Rating Confusion Matrix}
\begin{table}[H]\centering\scriptsize
\caption{XGBoost rating model confusion matrix on the repository split.}
\label{tab:rating-confusion}
\begin{tabular}{lrrrr}
\toprule
Actual / predicted & exceeds & exceptional & meets & needs improvement\\
\midrule
exceeds expectations & 525 & 12 & 7 & 0\\
exceptional & 5 & 443 & 0 & 0\\
meets expectations & 12 & 0 & 475 & 1\\
needs improvement & 0 & 0 & 7 & 513\\
\bottomrule
\end{tabular}
\end{table}

\section{Promotion Confusion Matrix}
\begin{table}[H]\centering\small
\caption{XGBoost promotion model confusion matrix on the repository split.}
\label{tab:promotion-confusion}
\begin{tabular}{lrr}
\toprule
Actual / predicted & not ready & ready\\
\midrule
not ready & 1005 & 245\\
ready & 259 & 491\\
\bottomrule
\end{tabular}
\end{table}

\section{Feature Importance}
\begin{table}[H]\centering\small
\caption{XGBoost rating feature importance from saved artifact.}
\label{tab:rating-feature-importance}
\begin{tabular}{lr}
\toprule
Feature & Importance\\
\midrule
\texttt{kpi\_score} & 0.6108\\
\texttt{task\_completion\_percent} & 0.1700\\
\texttt{goal\_completion\_percent} & 0.1085\\
\texttt{avg\_checkin\_progress} & 0.0535\\
\texttt{positive\_feedback\_ratio} & 0.0360\\
\texttt{tasks\_on\_time\_percent} & 0.0105\\
\texttt{checkin\_count} & 0.0061\\
\texttt{feedback\_count} & 0.0046\\
\bottomrule
\end{tabular}
\end{table}

\section{Interpretation}
The rating classifier performs very well because the synthetic target is strongly linked to the same input features that define the deterministic overall score. The promotion classifier performs more modestly because the dataset intentionally includes probabilistic promotion readiness. This is realistic as a demonstration, but it also means the service should be presented as an educational decision-support feature until real organizational validation exists.
